\documentclass[a4paper,12pt,twoside]{memoir}

% Castellano
\usepackage[spanish,es-tabla]{babel}
\usepackage{graphicx}
\usepackage{float}
\selectlanguage{spanish}
\usepackage[utf8]{inputenc}
\usepackage[T1]{fontenc}
\usepackage{lmodern} % Scalable font
\usepackage{microtype}
\usepackage{placeins}

\usepackage{url}

\RequirePackage{booktabs}
\RequirePackage[table]{xcolor}
\RequirePackage{xtab}
\RequirePackage{multirow}

% Links
\PassOptionsToPackage{hyphens}{url}\usepackage[colorlinks]{hyperref}
\hypersetup{
	allcolors = {red}
}

% Bibliography management
\usepackage[numbers,sort]{natbib}

% Ecuaciones
\usepackage{amsmath}

% Rutas de fichero / paquete
\newcommand{\ruta}[1]{{\sffamily #1}}

% Párrafos
\nonzeroparskip

% Huérfanas y viudas
\widowpenalty100000
\clubpenalty100000

% Imagenes
\usepackage{graphicx}
\newcommand{\imagen}[2]{
	\begin{figure}[!h]
		\centering
		\includegraphics[width=0.9\textwidth]{#1}
		\caption{#2}\label{fig:#1}
	\end{figure}
	\FloatBarrier
}

\newcommand{\imagenflotante}[2]{
	\begin{figure}%[!h]
		\centering
		\includegraphics[width=0.9\textwidth]{#1}
		\caption{#2}\label{fig:#1}
	\end{figure}
}



% El comando \figura nos permite insertar figuras comodamente, y utilizando
% siempre el mismo formato. Los parametros son:
% 1 -> Porcentaje del ancho de página que ocupará la figura (de 0 a 1)
% 2 --> Fichero de la imagen
% 3 --> Texto a pie de imagen
% 4 --> Etiqueta (label) para referencias
% 5 --> Opciones que queramos pasarle al \includegraphics
% 6 --> Opciones de posicionamiento a pasarle a \begin{figure}
\newcommand{\figuraConPosicion}[6]{%
  \setlength{\anchoFloat}{#1\textwidth}%
  \addtolength{\anchoFloat}{-4\fboxsep}%
  \setlength{\anchoFigura}{\anchoFloat}%
  \begin{figure}[#6]
    \begin{center}%
      \Ovalbox{%
        \begin{minipage}{\anchoFloat}%
          \begin{center}%
            \includegraphics[width=\anchoFigura,#5]{#2}%
            \caption{#3}%
            \label{#4}%
          \end{center}%
        \end{minipage}
      }%
    \end{center}%
  \end{figure}%
}

%
% Comando para incluir imágenes en formato apaisado (sin marco).
\newcommand{\figuraApaisadaSinMarco}[5]{%
  \begin{figure}%
    \begin{center}%
    \includegraphics[angle=90,height=#1\textheight,#5]{#2}%
    \caption{#3}%
    \label{#4}%
    \end{center}%
  \end{figure}%
}
% Para las tablas
\newcommand{\otoprule}{\midrule [\heavyrulewidth]}
%
% Nuevo comando para tablas pequeñas (menos de una página).
\newcommand{\tablaSmall}[5]{%
 \begin{table}
  \begin{center}
   \rowcolors {2}{gray!35}{}
   \begin{tabular}{#2}
    \toprule
    #4
    \otoprule
    #5
    \bottomrule
   \end{tabular}
   \caption{#1}
   \label{tabla:#3}
  \end{center}
 \end{table}
}

%
% Nuevo comando para tablas pequeñas (menos de una página).
\newcommand{\tablaSmallSinColores}[5]{%
 \begin{table}[H]
  \begin{center}
   \begin{tabular}{#2}
    \toprule
    #4
    \otoprule
    #5
    \bottomrule
   \end{tabular}
   \caption{#1}
   \label{tabla:#3}
  \end{center}
 \end{table}
}

\newcommand{\tablaApaisadaSmall}[5]{%
\begin{landscape}
  \begin{table}
   \begin{center}
    \rowcolors {2}{gray!35}{}
    \begin{tabular}{#2}
     \toprule
     #4
     \otoprule
     #5
     \bottomrule
    \end{tabular}
    \caption{#1}
    \label{tabla:#3}
   \end{center}
  \end{table}
\end{landscape}
}

%
% Nuevo comando para tablas grandes con cabecera y filas alternas coloreadas en gris.
\newcommand{\tabla}[6]{%
  \begin{center}
    \tablefirsthead{
      \toprule
      #5
      \otoprule
    }
    \tablehead{
      \multicolumn{#3}{l}{\small\sl continúa desde la página anterior}\\
      \toprule
      #5
      \otoprule
    }
    \tabletail{
      \hline
      \multicolumn{#3}{r}{\small\sl continúa en la página siguiente}\\
    }
    \tablelasttail{
      \hline
    }
    \bottomcaption{#1}
    \rowcolors {2}{gray!35}{}
    \begin{xtabular}{#2}
      #6
      \bottomrule
    \end{xtabular}
    \label{tabla:#4}
  \end{center}
}

%
% Nuevo comando para tablas grandes con cabecera.
\newcommand{\tablaSinColores}[6]{%
  \begin{center}
    \tablefirsthead{
      \toprule
      #5
      \otoprule
    }
    \tablehead{
      \multicolumn{#3}{l}{\small\sl continúa desde la página anterior}\\
      \toprule
      #5
      \otoprule
    }
    \tabletail{
      \hline
      \multicolumn{#3}{r}{\small\sl continúa en la página siguiente}\\
    }
    \tablelasttail{
      \hline
    }
    \bottomcaption{#1}
    \begin{xtabular}{#2}
      #6
      \bottomrule
    \end{xtabular}
    \label{tabla:#4}
  \end{center}
}

%
% Nuevo comando para tablas grandes sin cabecera.
\newcommand{\tablaSinCabecera}[5]{%
  \begin{center}
    \tablefirsthead{
      \toprule
    }
    \tablehead{
      \multicolumn{#3}{l}{\small\sl continúa desde la página anterior}\\
      \hline
    }
    \tabletail{
      \hline
      \multicolumn{#3}{r}{\small\sl continúa en la página siguiente}\\
    }
    \tablelasttail{
      \hline
    }
    \bottomcaption{#1}
  \begin{xtabular}{#2}
    #5
   \bottomrule
  \end{xtabular}
  \label{tabla:#4}
  \end{center}
}



\definecolor{cgoLight}{HTML}{EEEEEE}
\definecolor{cgoExtralight}{HTML}{FFFFFF}

%
% Nuevo comando para tablas grandes sin cabecera.
\newcommand{\tablaSinCabeceraConBandas}[5]{%
  \begin{center}
    \tablefirsthead{
      \toprule
    }
    \tablehead{
      \multicolumn{#3}{l}{\small\sl continúa desde la página anterior}\\
      \hline
    }
    \tabletail{
      \hline
      \multicolumn{#3}{r}{\small\sl continúa en la página siguiente}\\
    }
    \tablelasttail{
      \hline
    }
    \bottomcaption{#1}
    \rowcolors[]{1}{cgoExtralight}{cgoLight}

  \begin{xtabular}{#2}
    #5
   \bottomrule
  \end{xtabular}
  \label{tabla:#4}
  \end{center}
}



\graphicspath{ {./img/} }

% Capítulos
\chapterstyle{bianchi}
\newcommand{\capitulo}[2]{
	\setcounter{chapter}{#1}
	\setcounter{section}{0}
	\setcounter{figure}{0}
	\setcounter{table}{0}
	\chapter*{#2}
	\addcontentsline{toc}{chapter}{#2}
	\markboth{#2}{#2}
}

% Apéndices
\renewcommand{\appendixname}{Apéndice}
\renewcommand*\cftappendixname{\appendixname}

\newcommand{\apendice}[1]{
	%\renewcommand{\thechapter}{A}
	\chapter{#1}
}

\renewcommand*\cftappendixname{\appendixname\ }

% Formato de portada
\makeatletter
\usepackage{xcolor}
\newcommand{\tutor}[1]{\def\@tutor{#1}}
\newcommand{\course}[1]{\def\@course{#1}}
\definecolor{cpardoBox}{HTML}{E6E6FF}
\def\maketitle{
  \null
  \thispagestyle{empty}
  % Cabecera ----------------
\noindent\includegraphics[width=\textwidth]{cabecera}\vspace{1cm}%
  \vfill
  % Título proyecto y escudo informática ----------------
  \colorbox{cpardoBox}{%
    \begin{minipage}{.8\textwidth}
      \vspace{.5cm}\Large
      \begin{center}
      \textbf{TFG del Grado en Ingeniería Informática}\vspace{.6cm}\\
      \textbf{\LARGE\@title{}}
      \end{center}
      \vspace{.2cm}
    \end{minipage}

  }%
  \hfill\begin{minipage}{.20\textwidth}
    \includegraphics[width=\textwidth]{escudoInfor}
  \end{minipage}
  \vfill
  % Datos de alumno, curso y tutores ------------------
  \begin{center}%
  {%
    \noindent\LARGE
    Presentado por \@author{}\\ 
    en Universidad de Burgos --- \@date{}\\
    Tutores: \@tutor{}\\
  }%
  \end{center}%
  \null
  \cleardoublepage
  }
\makeatother

\newcommand{\nombre}{Roberto Arasti Blanco} %%% cambio de comando

% Datos de portada
\title{eLearningQA - Control de calidad}
\author{\nombre}
\tutor{Raúl Marticorena Sánchez y Carlos López Nozal}
\date{\today}

\begin{document}

\maketitle


\newpage\null\thispagestyle{empty}\newpage


%%%%%%%%%%%%%%%%%%%%%%%%%%%%%%%%%%%%%%%%%%%%%%%%%%%%%%%%%%%%%%%%%%%%%%%%%%%%%%%%%%%%%%%%
\thispagestyle{empty}


\noindent\includegraphics[width=\textwidth]{cabecera}\vspace{1cm}

\noindent D. Raúl Marticorena Sánchez, profesor del departamento de Ingeniería Informática,
área de Lenguajes y Sistemas Informáticos.

\noindent D. Carlos López Nozal, profesor del departamento de Ingeniería Informática,
área de Lenguajes y Sistemas Informáticos.

\noindent Exponen:

\noindent Que el alumno D. \nombre, con DNI 71307060E, ha realizado el Trabajo final de Grado en Ingeniería Informática titulado \textit{eLearningQA - Control de calidad}. 

\noindent Y que dicho trabajo ha sido realizado por el alumno bajo la dirección del que suscribe, en virtud de lo cual se autoriza su presentación y defensa.

\begin{center} %\large
En Burgos, {\large \today}
\end{center}

\vfill\vfill\vfill

% Author and supervisor
\begin{minipage}{0.45\textwidth}
\begin{flushleft} %\large
Vº. Bº. del Tutor:\\[2cm]
D. Raúl Marticorena Sánchez
\end{flushleft}
\end{minipage}
\hfill
\begin{minipage}{0.45\textwidth}
\begin{flushleft} %\large
Vº. Bº. del tutor:\\[2cm]
D. Carlos López Nozal
\end{flushleft}
\end{minipage}
\hfill

\vfill

% para casos con solo un tutor comentar lo anterior
% y descomentar lo siguiente
%Vº. Bº. del Tutor:\\[2cm]
%D. nombre tutor


\newpage\null\thispagestyle{empty}\newpage




\frontmatter

% Abstract en castellano
\renewcommand*\abstractname{Resumen}
\begin{abstract}
El proceso de enseñanza y aprendizaje en línea, conocido como e-learning, se ha implantado de manera acelerada en muchas instituciones académicas, debido en gran medida por la situación de confinamiento provocada por la pandemia COVID-19. Las instituciones educativas tradicionales se han visto obligadas a impartir sus cursos en plataformas virtuales de aprendizaje, como Moodle o similares. La calidad del e-learning es un aspecto fundamental para la enseñanza hoy en día. Existen frameworks para valorar la calidad de un curso en línea desde distintas perspectivas (pedagógica, tecnológica y estratégica), identificando roles en el curso (diseñador, facilitador y proveedor) y en las diferentes fases (análisis, diseño, implementación y evaluación). A pesar de esto, el control de la calidad en este ámbito todavía no cuenta con herramientas capaces de ayudar a la toma de decisiones de forma automática. Bajo esta necesidad de valorar la calidad de los cursos en línea y marco teórico de calidad, este proyecto tiene como objetivo principal adaptar un framework de calidad definiendo métricas obtenidas con consultas sobre los cursos de un LMS. Además se desarrolla una aplicación software que dado un curso o cursos genere una evaluación de calidad por medio de servicios web y otras herramientas si es necesario.
\end{abstract}

\renewcommand*\abstractname{Descriptores}
\begin{abstract}
Aplicación web, framework de calidad en e-learning, revisiones automáticas, calidad en cursos en línea, Learning Management System (LMS).
\end{abstract}

\clearpage

% Abstract en inglés
\renewcommand*\abstractname{Abstract}
\begin{abstract}
The online teaching and learning process, known as e-learning, has been implanted in a hurrisome way in many academic institutions, due in great measure to the situation of confinement triggered by the COVID-19 pandemic. Traditional educational institutions have been forced to impart their courses on virtual learning platforms, such as Moodle o similar. The quality of e-learning is a fundamental aspect of teaching nowadays. There are frameworks to assess the quality of an online course from different perspectives (pedagogical, technological and strategic), identifying roles in the course (designer, facilitator and provider) and in the different phases (analysis, design, implementation and evaluation). Despite this, quality assurance in this area is unprovided of tools able to help in the decision making process in an automatic way. Under this necessity of assessing online course quality and theoretical framework of quality, this project has as a main objective to adapt a quality framework defining metrics obtained with queries about courses in an LMS. In addition, a software application is developed that, given a course or courses, generates a quality evaluation via web services and other tools if needed.
\end{abstract}

\renewcommand*\abstractname{Keywords}
\begin{abstract}
Web application, e-learning quality framework, automatic revisions, online course quality, Learning Management System (LMS).
\end{abstract}

\clearpage

% Indices
\tableofcontents

\clearpage

\listoffigures

\clearpage

\listoftables
\clearpage

\mainmatter
\capitulo{1}{Introducción}

Este trabajo consiste en una aplicación de navegador que a partir de una serie de métricas analice la calidad del diseño instruccional en las distintas asignaturas de un profesor en Moodle y pueda generar informes.

\capitulo{2}{Objetivos del proyecto}

El objetivo principal de este trabajo es desarrollar una aplicación de navegador
que permita al profesor evaluar su propio diseño instruccional a través de sus asignaturas en Moodle.
Un segundo objetivo es una aplicación aparte para la visualización de informes.

Los objetivos técnicos de este trabajo para conseguir los objetivos son:
\begin{itemize}
	\item Ser capaces de extraer información de la web para analizarla.
	\item Poder generar informes de forma automática.
\end{itemize}

\capitulo{3}{Conceptos teóricos}


\section{Definiciones}

\subsection{E-learning}

El e-learning es la enseñanza impartida por medios electrónicos como internet, plataformas virtuales, medios audiovisuales…etc.

\subsection{Moodle}

Moodle es una plataforma de aprendizaje que permite a los profesores crear entornos de aprendizaje altamente personalizables. Fue creado por Martin Dougiamas que publicó su primera versión el 20 de agosto de 2002\cite{dougiamas2002interpretive}.

\subsection{Calidad}

La calidad se puede definir como la capacidad de satisfacer una serie de necesidades y en el caso del e-learning se trata de las necesidades educativas del alumno, como la calidad del material educativo o la ayuda a la comprensión de este.

\subsection{Web scraping}

Es el conjunto de técnicas utilizadas para extraer y almacenar información de la web, un programa que hace esto de forma automática se llama web crawler.

\subsection{Web services}

Es un medio de comunicación entre dos aplicaciones en ordenadores distintos dentro de una red mediante el uso de distintos protocolos. Se puede pedir cierta información o acciones al servidor por medio de llamadas a funciones. En el caso de los web services que proporciona Moodle existe un servidor que hace de intermediario entre cliente y proveedor\cite{moodle-2020}.

Ahora paso a explicar las tres dimensiones del marco de referencia de calidad del e-learning del MOOQ\cite{stracke2018quality}, en el que me baso:
\subsection{Fases}

\subsubsection{Análisis}
En esta fase se definen los objetivos, el contexto, y los recursos(docentes, tiempo, presupuesto...) para la ejecución para comprender la situación inicial.

\subsubsection{Diseño}
En esta fase se define lo que se planea hacer a partir de los resultados de la fase de análisis, como por ejemplo el enfoque de enseñanza que se piensa llevar a cabo.

\subsubsection{Implementación}
En esta fase se define de qué manera se van a llevar a cabo los planes descritos en la fase de diseño, como por ejemplo de qué manera se va a producir el contenido.

\subsubsection{Realización}
Esta fase es en la que se interactúa con el alumno, se gestionan los problemas técnicos y las dudas de los alumnos, además de evaluar su aprendizaje.

\subsubsection{Evaluación}
En esta última fase se evalúa la calidad del resto de fases mediante encuestas, entrevistas, u otros medios.

\subsection{Roles}

\subsubsection{Diseñadores}
Los encargados de decidir de qué forma se van a impartir el curso y generan el contenido(Autores, expertos en el tema, diseñadores instruccionales). 

\subsubsection{Facilitadores}
Son aquellos que conocen la materia a enseñar y son capaces de explicarlo y dar feedback, además de seguir el aprendizaje de los alumnos.

\subsubsection{Proveedores}
Son los encargados de proveer los medios digitales para llevar a cabo el e-learning(programadores, diseñadores y desarrolladores de software).

\subsection{Perspectivas}

\subsubsection{Pedagógica}
El punto de vista que se centra en el contenido y en el aprendizaje por parte del alumno. Los procesos relacionados tienen que ver con el contenido, el feedback, y la evaluación de los alumnos.

\subsubsection{Tecnológica}
El punto de vista que se centra en las necesidades tecnológicas del curso. La mayoría de los procesos está relacionada con esta perspectiva debido a la naturaleza del e-learning.

\subsubsection{Estratégica}
El punto de vista que se centra en la consecución de los objetivos del curso dentro del tiempo y presupuesto establecidos para este. Los procesos relacionados tienen que ver con los objetivos, conceptos y el contexto en el que se enseña(presupuesto, demanda, competencia...).


\section{Consultas}
En este apartado definiré qué información intentamos conseguir de cada fase, rol y perspectiva mediante distintos métodos.
Leyenda:
Responsabilidad:R=Responsable,I=Involucrado

\subsection{Analisis:} 
Consulta:La división en grupos es correcta 

Queries:Comprobar en los logs que no hay ningún alumno que no pertenezca a ningún grupo

Perspectivas:Pedagógica, Estratégica 
Diseñador:I Facilitador:\_ Proveedor:R

\subsection{Diseño:}
Consulta:Las opciones de progreso del estudiante están activadas

Queries:Llamar a la función core\_completion\_get\_activities\_completion\_status

Perspectivas:Pedagógica
Diseñador:R Facilitador:I Proveedor:\_
\subsubsection{}
Consulta:Se proporcionan contenidos en diferentes formatos

Queries:Comprobar mediante webscraping que hay contenido audiovisual embedido

Perspectivas:Pedagógica, Tecnológica 
Diseñador:R Facilitador:I Proveedor:I

\subsection{Implementación:}
Consulta:Los recursos estan actualizados

Queries:Accede a las fechas de los archivos para comprobar si son recientes o no

Perspectivas:Pedagógica, Tecnológica 
Diseñador:R Facilitador:I Proveedor:I

\subsection{Realización:}
Consulta:Fechas de apertura y cierre de tareas son correctas

Queries:Comprobar que las fechas son de hace menos de 6 meses y que la de apertura es anterior a la de cierre

Perspectivas:Pedagógica, Tecnológica 
Diseñador:I Facilitador:R Proveedor:I
\subsubsection{}
Consulta:Se detallan los criterios de evaluación (rúbricas, ejemplos)

Queries:Comprobar mediante webscraping que hay una rúbrica en la tarea

Perspectivas:Pedagógica, Tecnológica 
Diseñador:R Facilitador:I Proveedor:I
\subsubsection{}
Consulta:El profesor responde en los foros dentro del límite de 48 horas lectivas desde que se plantea la duda

Queries:Comprobación de la actividad en los foros en los logs

Perspectivas:Pedagógica, Tecnológica 
Diseñador:I Facilitador:R Proveedor:\_
\subsubsection{}
Consulta:Se ofrece retroalimentación de las tareas

Queries:Comprobación en los logs de que haya un comentario del profesor en las tareas cerradas hace una semana

Perspectivas:Pedagógica, Tecnológica 
Diseñador:\_ Facilitador:R Proveedor:\_
\subsubsection{}
Consulta:Las tareas están calificadas

Queries:Comprobación en los logs de que se hayan calificado todas las tareas cerradas hace una semana

Perspectivas:Pedagógica, Tecnológica 
Diseñador:\_ Facilitador:R Proveedor:\_

\subsection{Evaluación:}
Consulta:La mayoría de alumnos responden a los feedbacks

Queries:Llamar a la función mod\_feedback\_get\_non\_respondents

Perspectivas:Pedagógica, Tecnológica, Estratégica
Diseñador:I Facilitador:R Proveedor:I



\capitulo{4}{Técnicas y herramientas}
\section{Técnicas}
He dividido el trabajo en sprints de una duración de 14 días cada uno, teniendo reuniones con los tutores entre sprint y sprint. También he implementado integración continua y despliegue continuo en el proyecto.
En el desarrollo del proyecto he priorizado el funcionamiento del software sobre la documentación extensiva y la respuesta al cambio sobre ceñirse a un plan. Todos estos elementos se pueden asociar a las metodologías ágiles, sin embargo, no opino que el desarrollo de este proyecto haya seguido una metodología específica.
\section{Entorno de desarrollo integrado: IntelliJ IDEA}
Para el desarrollo de la aplicación web en Java, he elegido IDEA en oposición a Eclipse.
A pesar de que conozco Eclipse desde hace tres años y esta es la primera vez que utilizo IDEA, el año pasado utilicé Android Studio para desarrollar una aplicación Android, durante ese desarrollo me dí cuenta de que estaba basado en IDEA y que me sentía mucho más cómodo programando en Java en IDEs de JetBrains que en Eclipse. Facilita bastante la refactorización y escritura del código debido a funciones como el completado de código y las acciones rápidas, funcionalidades que también posee Eclipse, pero que en IDEA son más intuitivas y satisfactorias.
\section{Patrón de diseño: Fachada}
El patrón de diseño fachada consiste en crear una clase que haga de intermediario entre el cliente y uno o varios subsistemas de la aplicación con varios propósitos: Simplificar y centralizar el control, actuar como elemento de seguridad restringiendo el acceso a ciertas partes, y separar responsabilidades de los subsistemas. Un mismo sistema podría tener varias fachadas distintas que den un mismo servicio de distintas formas, por ejemplo, mi fachada se utiliza para generar parte del contenido de mi aplicación web, pero si quisiera trasladar la aplicación a una de escritorio solo tendría que crear una nueva fachada dejando los sistemas subyacentes intactos.
\section{Hospedaje del repositorio: GitHub}
Github es un sitio de hospedaje de repositorios de código gratis para repositorios públicos.
GitHub Actions, integrado en GitHub, premite definir flujos de trabajo y lo he utilizado para implementar la integración continua.
\section{Control de versiones: GitHub Desktop}
GitHub Desktop es una aplicación que proporciona una interfaz gráfica para interactuar con tus repositorios en GitHub. Simplifica las acciones más comunes asociándoles botones, muestra todas las diferencias introducidas antes de realizar un commit, gestiona tus credenciales de forma automática, y permite alternar entre ramas y repositorios en dos clicks, con lo que se agiliza mucho el control de versiones.
\section{Hospedaje de la aplicación: Heroku}
Heroku es una plataforma de hospedaje y computación de aplicaciones web dinámicas. Tiene una versión para hospedar una aplicación no comercial y poder ejecutarla durante 500 horas al mes de forma gratuita.
\section{Ejecución de tests: JUnit}
JUnit es un framework para la realización de tests en Java. Es compatible con varios entornos de desarrollo integrados e incluso se puede usar mediante linea de comandos.
\section{Construcción automática del proyecto: Maven}
Maven es una herramienta para la construcción de proyectos software. Utiliza un archivo definido dentro del proyecto llamado ''pom.xml'' para definir la configuración necesaria para construir el proyecto como las dependencias o el formato al que compilar.
\section{Generación de tablas: TablesGenerator.com}
TablesGenerator.com es una página que permite crear una tabla con facilidad y puede convertirla a distintos formatos (LaTeX, HTML, texto simple...). La he utilizado para diseñar las tablas utilizadas en la aplicación y las tablas en memoria y anexos.
\section{Paso de JSON a POJO: Json2CSharp.com}
Json2CSharp.com es una página que permite la conversión de un objeto JSON a su POJO en Java correspondiente entre otras. La he utilizado en el proceso de implementación de las comprobaciones sobre los cursos Moodle debido a que las respuestas a mis llamadas REST tenían que ser deserializadas a objetos para poder manejarlas con facilidad.
\section{Redacción en HTML: HtmlNotepad}
HtmlNotepad es un editor HTML que facilita la redacción de textos en HTML. Contiene gran cantidad de atajos de teclado que agiliza mucho el formateo de texto. Lo he utilizado para la redacción de la página del manual de usuario.
\section{Generación de diagramas UML: ArgoUML}
ArgoUML es una herramienta de dibujo de diagramas UML. Lo he utilizado para todos los diagramas de los anexos. 
\section{Redacción de memoria y anexos: TeXstudio}
TeXstudio es un editor de LaTeX. Contiene cantidad de herramientas necesarias para la creación de documentos en LaTeX y cuenta con un visor de PDF.
 
 
 
 
 
 
\section{Herramientas para acceder a la información}
\begin{itemize}
	\item \textbf{Web scraping:}
	Es el conjunto de técnicas utilizadas para extraer y almacenar información de la web, un programa que hace esto de forma automática se llama web crawler. Algunas de estas técnicas nos permiten obtener información que no podríamos sacar con otras herramientas como por ejemplo si una página contiene enlaces fuera de Moodle a partir de su código HTML.
	\item \textbf{Web services:}
	Es un medio de comunicación entre dos aplicaciones en ordenadores distintos dentro de una red mediante el uso de distintos protocolos. Se puede pedir cierta información o acciones al servidor por medio de llamadas a funciones. En el caso de los web services que proporciona Moodle existe un servidor que hace de intermediario entre cliente y proveedor \cite{moodle-2020}. Por ejemplo, llamando a la función core\_grades\_get\_grades podríamos obtener las notas de un alumno.
	\item \textbf{Logs:}
	Son registros de actividad que Moodle crea a partir de los eventos que realizan los usuarios, como publicar en un foro o empezar un cuestionario. Se pueden descargar en forma de archivo. Un ejemplo del uso que le daríamos sería obtener la actividad de un profesor para saber si responde a las dudas de los alumnos en los foros.
\end{itemize}
\section{Framework de desarrollo web}
Respecto a qué framework utilizar para el desarrollo de la aplicación web hemos barajado estas opciones:
\begin{itemize}
	\item \textbf{Spring:}
	Es un framework para la creación de aplicaciones Java. Gestiona las dependencias entre objetos de forma automática lo que permite un bajo nivel de acoplamiento y ofrece un framework para desarrollo de aplicaciones web. Su fuerte es la escalabillidad, pero es complicado aprender a usarlo.
	\item \textbf{Vaadin:}
	Es un framework de desarrollo web con una gran librería de componentes web. Usa GWT (Google Web Toolkit) para compilar Java a JavaScript y evitar al programador usar otros lenguajes aparte de Java. Su versión de pago ofrece una herramienta de edición gráfica que acelera el proceso de creación. Su fuerte es la velocidad de creación de prototipos, pero tiene mala escalabilidad.
\end{itemize}

De entre estas dos me decanto por Spring, ya que la escalabilidad es bastante mejor y es mucho más popular que Vaadin, lo que me ayudará a encontrar documentación y tutoriales a la hora de encontrar problemas durante el desarrollo.
\capitulo{5}{Aspectos relevantes del desarrollo del proyecto}
\section{Ciclo de vida}
La realización de este trabajo se ha llevado a cabo en sprints de una duración de 14 días con reuniones entre sprint y sprint, y con frecuencia han habido reuniones a mitad de estos.

El proyecto empezó sintetizando una lista de aspectos a comprobar en los cursos Moodle a partir del marco de referencia de calidad del e-learning de MOOQ \cite{stracke2018quality} y un documento interno de UBUCEV proporcionado por los tutores.
Después tuve que decidir si crear una aplicación de escritorio o una aplicación web y qué framework utilizar para desarrollar la aplicación.
A partir de ahí, se hizo un prototipo para comprobar que era capaz de acceder a los web services de Moodle desde mi aplicación. Este prototipo solo mostraba la lista de cursos accedidos recientemente a partir de unas credenciales para la página de demostración de Moodle llamada Mount Orange School. Más tarde establecí el ciclo de integración continua/despliegue continuo y lancé la aplicación en Heroku. Una semana después, al final del sprint había creado una versión del informe específico que solo realizaba una comprobación sobre los cursos. A partir de ahí, debido a las fechas (segunda quincena de diciembre) no hubo ninguna reunión con los tutores hasta la vuelta de las vacaciones de navidad. Durante ese periodo, añadí el resto de las comprobaciones al informe y me dediqué a completar partes de la memoria.

\section{Proceso de obtención de llamadas a los servicios web}
Para obtener la información necesaria para las comprobaciones sobre los cursos de Moodle he tenido que realizar llamadas REST a distintas funciones de la API de servicios web de Moodle \cite{wsapifunctions-2021}. Sin embargo, la tabla que detalla la lista de funciones de la API omite los parámetros necesarios para las llamadas a las funciones.
Para averiguar qué parametros debía utilizar en cada función tuve que acceder al repositorio de Moodle \cite{moodlerepository-2022} y encontrar las funciones que detallaban el funcionamiento de la función de la API en la que estaba interesado cada vez.
El nombre de una función se puede dividir en dos partes: el nombre del componente que posee la función, y el nombre de la función. Las funciones están escritas en PHP y aparecen dentro de los archivos llamados "externallib.php". Para cada función existen tres funciones asociadas: <NOMBRE DE LA FUNCIÓN>, <NOMBRE DE LA FUNCIÓN>\_parameters, y <NOMBRE DE LA FUNCIÓN>\_returns.
La función a secas define el comportamiento de la función, y en su declaración se puede ver qué parámetros son opcionales y cuáles no.
La función con ``parameters'' al final describe los nombres y los tipos de datos de los parámetros que espera recibir la función.
La función con ``returns'' al final describe los nombres y los tipos de datos de los atributos que devuelve la función.
Por ejemplo, si quiero averiguar como llamar a la función mod\_forum\_get\_forums\_by\_courses
tengo acceder al fichero /mod/forum/externallib.php (supuesto por el nombre del componente que contiene la función, ``mod\_forum'' en este caso) y buscar las funciones get\_forums\_by\_courses, get\_forums\_by\_courses\_parameters, y get\_forums\_by\_courses\_returns.

Para probar el funcionamiento de las llamadas y al mismo tiempo obtener las clases que debía definir como POJO para manejar los datos recibidos como JSON (JSON (JavaScript Object Notation), es una forma de representar un objeto en formato de texto muy usado para transmitir información en aplicaciones web), hice lo siguiente:
Uso la página de demostración de Moodle llamada Mount Orange School para hacer pruebas.
Hago llamadas REST de forma manual con el navegador Chrome y obtengo el token con las credenciales de profesor (usuario:``teacher'' contraseña:``moodle'')\imagen{Token.png}{Obtención del token}. Después utilizo el token para la siguiente llamada y obtengo una respuesta en formato JSON. \imagen{JSON.png}{Obtención del JSON} Copio la respuesta en un conversor de JSON a POJO.\imagen{Conversor.png}{Uso del conversor} Creo las clases correspondientes, añado un constructor vacío y encapsulo los atributos. 


\section{Integración continua y despliegue continuo}
La integración continua consiste en la automatización de la compilación y ejecución de pruebas cada vez que se suben cambios al repositorio.
El despliegue continuo es la automatización del despliegue de un producto tras cada cambio en el repositorio.
He implementado la integración continua del proyecto con GitHub Actions, primero, establecí en el archivo pom.xml que la compilación del proyecto fuera en formato WAR (Web Application Resource), luego, creé el archivo "maven.yml" en la carpeta de workflows para establecer que cada vez que se realice un push en la rama principal el proyecto se compile y se ejecuten los tests con Maven.
He implementado el despliegue continuo en Heroku, la mayoría del proceso ha sido bastante intuitiva, ya que una de las opciones que ofrecía era GitHub como método de despliegue pero debido a que mi repositorio contiene por un lado la memoria y por otro el proyecto software, he tenido que añadir un buildpack (conjuntos de scripts de codigo abierto usados para compilar las aplicaciones en Heroku) que permite especificar una subcarpeta del repositorio para usarla como directorio raíz del proyecto software. También establecí en las opciones que el despliegue automático espere a que se supere la integración continua.

\section{Uso del nombre Moodle}
Como el objetivo inicial del proyecto es integrar Moodle en la aplicación, el nombre de tanto el proyecto como la aplicación iba a ser ``MoodleQA'', pero en medio del desarrollo los tutores y yo descubrimos que al ser ``Moodle'' una marca registrada \cite{moodletrademark-2022}, existe cierta cantidad de restricciones sobre el uso de esta palabra. Las restricciones que nos conciernen son no poder utilizar ``Moodle'' en el nombre de tu software ni en el nombre de dominio ni usarlo para describir tu aplicación de forma que los usuarios crean que estás asociado con Moodle si no lo estás. Por ello he tenido que revisar todas las menciones que hago a Moodle tanto en la memoria y anexos como en la aplicación.

\section{Informes: pestañas vs breadcrumbs}
A la hora de mostrar los informes se presenta la disyuntiva de si hacerlo en una sola pestaña de navegador y permitir al usuario navegar entre los distintos cursos con ayuda de una miga de pan o breadcrumb, o por el contrario generar cada informe en una nueva pestaña. Me he decantado por generar una pestaña por informe por dos razones: el rendimiento y la comparación de informes. Cada vez que se genera un informe se descarga la información necesaria para las consultas y se realizan las consultas sobre dicha información, para un servidor puede que la carga no sea mucha y se puedan generar los informes al instante, pero ejecutando la aplicación en local desde mi ordenador portátil he llegado a tener esperas de alrededor de diez segundos por un solo informe. A la hora de comparar informes es más ágil el cambiar de pestaña que retroceder a la página anterior e ir al otro informe.

\section{Definición del quality gate de SonarCloud}
En SonarCloud la quality gate es una serie de condiciones sobre la calidad del codigo, por ejemplo, que la puntuación de seguridad no esté por debajo de ``A''. Existe una quality gate predefinida llamada Sonar way, pero está definida para proyectos reales, así que en mi caso he decidido definir mi propia quality gate a partir de otros proyectos en Sonarcloud de la misma índole. Realicé una búsqueda y encontré 25 TFGs de la UBU, y calculé las medianas de dichos trabajos para obtener las condiciones de mi quality gate.

\section{SonarCloud en el ciclo de desarrollo}
Desde que añadí SonarCloud al ciclo de integración y despliegue continuo he creado el hábito de comprobar los bugs, smells, y vulnerabilidades generados tras cada commit, esto evita que la cantidad de problemas se acumule y a su vez hace que el código sea más fácil de mantener.

\section{Diseño del logo de la aplicación}
Debido a que la aplicación realiza una serie de consultas se me ocurrió que el logo lo formasen unas aspas y un check, pero sin el uso de colores el símbolo resultante puede resultar confuso, así que rodeé la parte del check con color verde y la parte de las aspas con color rojo para hacerlo más entendible.
Los colores usados son los de Bootstrap para crear concordancia con el resto de la página y de paso crear las imágenes de check y aspas usados en los resultados de las consultas y que de esta forma se mimeticen en el color de fondo de la celda. El texto del logo es blanco con reborde negro para hacerlo legible con independencia del color del fondo.

\section{Almacenamiento de registros de informes}
La aplicación almacena los registros de los informes generados para poder generar los gráficos de evolución. Ejecutando la aplicación en local no supone ningún problema, sin embargo, el almacenamiento en Heroku es volátil \cite{ephemeralHeroku-2018}. La solución a esto sería contratar un proveedor de active storage como Amazon S3 para tener un almacenamiento que funcione independiente a la aplicación \cite{activeStorageOverview-2019}.

\capitulo{6}{Trabajos relacionados}

Este apartado resume los estudios y proyectos realizados por terceros que están relacionados con la garantía y control de calidad de los procesos de enseñanza y aprendizaje en cursos en línea, y con sus consecuencias. También se habla de las novedades que trae este proyecto a este ámbito.

\subsection{Automated e-learning quality evaluation}
Un articulo presentado por Rositsa Doneva y Silvia Gaftandzhieva en la conferencia internacional del e-learning celebrada en Berlin en septiembre de 2015.
En este artículo se llevan a cabo dos experimentos para probar la factibilidad de implementar un sistema automático de evaluación de la calidad en Moodle. El primer experimento consiste en integrar UBIS-Jaspersoft, un sistema de business intelligence para universidades, con la base de datos de Moodle para analizar los resultados de las respuestas de los alumnos en un modulo de feedback de Moodle. El segundo experimento consiste en la creación de cuatro servicios web para ser utilizados en la evaluación de distintos indicadores de la calidad \cite{doneva2015automated}.

\subsection{Perceived Service Quality and Student Loyalty in an Online University}
Un artículo presentado por María-Jesús Martínez-Argüelles y Josep-Maria Batalla-Busquets en la revista IRRODL en 2016.
Este artículo estudia la relación entre todos los aspectos de la enseñanza (incluida la interfaz de usuario en el e-learning) y la percepción de calidad del servicio por parte del estudiante, y entre esta última y la lealtad y la disposición a la recomendación por parte de este. El estudio llega a la conclusión de que hay una relación directa y no solo indirecta entre estas variables \cite{martinez2016perceived}.

\subsection{Dashboard for Evaluating the Quality of Open Learning Courses}
Un artículo presentado por Gina Mejía-Madrid, Faraón Llorens-Largo, y Rafael Molina-Carmona en la revista Sustainability en 2020.
Este artículo presenta un modelo para la evaluación de la calidad de los cursos de Open Learning y un \textit{dashboard} creado a partir de resultados de encuestas y entrevistas. La palabra ``dashboard'' se podría traducir como el panel de instrumentos de un coche u otro vehiculo, la analogía viene de qué el tipo de \textit{dashboard} al que nos estamos refiriendo es un conjunto de tablas y gráficos fáciles de leer que permiten a aquel que lo mira entender la información de forma rápida y tomar decisiones. En el artículo también se habla de forma muy breve de automatizar la obtención de datos para generar el \textit{dashboard} pero sin llegar a definir si pretenden automatizar las encuestas o obtener los datos de forma automática por otros medios \cite{mejia2020dashboard}.

\subsection{A Hierarchical Model to Evaluate the Quality of Web-Based E-Learning Systems}
Un artículo presentado por Muhammad Abdul Hafeez y otros seis autores en la revista Sustainability en 2020.
Este artículo presenta un modelo para definir la calidad de los sistemas de e-learning generado a partir de una serie de encuestas que les permitieron identificar los factores clave para la calidad según su importancia. El resultado es un modelo con forma de árbol en el que los nodos hoja son los aspectos a evaluar y que se encuentran ordenados por su relevancia dentro de su nodo padre \cite{muhammad2020hierarchical}.

\subsection{Data Analysis for Evaluation on Course Design and Improvement of “Cyberethics” Moodle Online Courses}
Un artículo presentado por Motonori Nakamura y Hiroshi Ueda en la revista Procedia Computer Science en 2017.
En este artículo se crea un sistema de recolección de datos con el objetivo de analizar los resultados en la recepción por parte de los estudiantes en un curso Moodle de ciberética de Japón para poder ver los efectos de los cambios realizados en el diseño del curso a lo largo de los años. El artículo concluye que los cambios conllevaron resultados concretos \cite{ueda2017data}.


\subsection{Moodle Course Checker Plugin}
Un plugin para Moodle que permite realizar una serie de comprobaciones mediante un conjunto de comprobadores independientes que se pueden ejecutar de forma individual \cite{coursechecker-2021}.

\subsection{Course Checks Block Plugin}
Otro plugin para Moodle que permite realizar una serie de comprobaciones automáticas \cite{coursechecksblock-2018}.

\subsection{Comparativa de herramientas relacionadas}
Este proyecto adapta un marco de calidad de cursos de e-learning a una serie de comprobaciones sobre cursos Moodle y es capaz de generar informes que además de mostrar los resultados de dichas comprobaciones. Muestra la traslación de estos a las responsabilidades de cada fase, rol, y perspectiva descritos en el marco y señala qué elementos se pueden cambiar para mejorar los resultados.

\begin{table}[H]
\resizebox{\textwidth}{!}{%
	\begin{tabular}{l|ccc}
		\hline
		\rowcolor[HTML]{FFFFFF} 
		\textbf{Característica}                                                                 & \textbf{eLearningQA} & \textbf{Course Checker}                                                        & \textbf{\begin{tabular}[c]{@{}c@{}}Course Checks\\ Block\end{tabular}}         \\ \hline
		\rowcolor[HTML]{EFEFEF} 
		Idiomas                                                                                 & Español           & \begin{tabular}[c]{@{}c@{}}Español, inglés,\\ alemán, y portugués\end{tabular} & \begin{tabular}[c]{@{}c@{}}Español, inglés,\\ portugués, y griego\end{tabular} \\
		Nº de comprobaciones                                                                    & 17                & 10                                                                             & 7                                                                              \\
		\rowcolor[HTML]{EFEFEF} 
		\begin{tabular}[c]{@{}l@{}}Ejecución independiente\\ de comprobaciones\end{tabular}     & No                & Sí                                                                             & No                                                                             \\
		\begin{tabular}[c]{@{}l@{}}Contiene enlaces para\\ solventar los problemas\end{tabular} & Sí                & Sí                                                                             & No                                                                             \\
		\rowcolor[HTML]{EFEFEF} 
		\begin{tabular}[c]{@{}l@{}}Versiones de Moodle\\ compatibles\end{tabular}               & v3.8+             & v3.6+                                                                          & v2.6+                                                                          \\
		Tipo                                                                                    & Aplicación web    & Plugin                                                                         & Plugin                                                                         \\ \hline
	\end{tabular}%
}
\end{table}
\capitulo{7}{Conclusiones y Líneas de trabajo futuras}

\section{Líneas de trabajo futuras}
\begin{itemize}
	\item
	El método que utiliza la aplicación para realizar la comprobación de si el profesor responde a las dudas de los alumnos es una solución preliminar e incompleta. Se tiene como objetivo a futuro encontrar una forma más fiable de determinar qué es una duda y cuándo ha sido resuelta.
	El uso de modelos basados en el procesamiento del lenguaje natural puede ser un campo exploratorio que permita poder clasificar un mensaje del foro como una respuesta de dudas de un profesor. Pensamos que el diseño experimental y el posterior análisis de un clasificador con este cometido es suficientemente complejo para ser considerado un TFG por sí mismo.
	\item
	De las cinco consultas definidas referentes a los cuestionarios Moodle, no hay ninguna implementada en la aplicación, y esto se debe a que al menos cuatro de ellas no están soportadas por la API de servicios web de Moodle.
	Una posible solución a estas comprobaciones es el uso de técnicas de web scraping, pero teniendo en cuenta que las páginas que contengan la información deseada pueden variar en función del servidor que aloje la plataforma, no es una solución trivial.
	\item
	Por el momento la aplicación solo integra cursos Moodle, pero sería conveniente que la aplicación permitiera analizar cursos online de otros LMS como Blackboard o Edmodo. Sin embargo, realizar los cambios para esto supondría adaptarse a las APIs de servicios correspondientes suponiendo que sean lo suficientemente parecidas, y en caso de no poder acceder a la información necesaria mediante servicios web, habría que implementar otras formas de acceder a la información necesaria.
	\item
	Los plugin de Moodle con los que comparo la aplicación tienen más idiomas disponibles, esto se debe a que están dispuestos de forma que cualquiera pueda aportar sus propias traducciones, sin embargo, el internacionalizar la aplicación la haría más competitiva.
	\item
	La lista de cursos que muestra la aplicación es la lista de los cursos en los que se encuentra matriculado el usuario con independencia de su rol o los permisos que tenga. Sería interesante seguir mostrando los mismos cursos pero sin resaltar en caso de que no se tengan los permisos necesarios para realizar las consultas para poder contactar al administrador en caso de problemas con los permisos.
	\item
	Sería muy interesante poder activar y desactivar las diferentes consultas desde los archivos de configuración si tenemos en cuenta que algunas consultas no aplican para algunos tipos de curso. Sin embargo, a la hora de calcular las estadísticas posteriores habría que adaptarse según las consultas que estén activadas, algo que es difícil teniendo en cuenta que la matriz Rol-Responsabilidad del informe se obtiene multiplicando un vector con los puntos obtenidos de cada consulta por una matriz que contiene la cantidad de puntos que supone el cumplimiento de cada comprobación para cada combinación de rol y perspectiva. El objetivo sería buscar otra forma de realizar el cálculo para facilitar la implementación de lo primero.
\end{itemize}


%\bibliographystyle{plain}
\bibliographystyle{plainurl}
\bibliography{bibliografia}

\end{document}
